\documentclass[10pt,a4paper]{amsart}
\usepackage{amsmath,amssymb,amsthm,graphicx}
\usepackage{verbatim}
\newenvironment{reference}{\comment}{\endcomment}
%\newenvironment{reference}{}{}
\newenvironment{slogan}{\comment}{\endcomment}
\newenvironment{history}{\comment}{\endcomment}

% For commutative diagrams we use Xy-pic
\usepackage[all]{xy}

% We use 2cell for 2-commutative diagrams.
\xyoption{2cell}
% Theorem environments.
%
\theoremstyle{plain}
\newtheorem{theorem}[subsection]{Theorem}
\newtheorem{proposition}[subsection]{Proposition}
\newtheorem{lemma}[subsection]{Lemma}

\theoremstyle{definition}
\newtheorem{definition}[subsection]{Definition}
\newtheorem{example}[subsection]{Example}
\newtheorem{exercise}[subsection]{Exercise}
\newtheorem{situation}[subsection]{Situation}

\theoremstyle{remark}
\newtheorem{remark}[subsection]{Remark}
\newtheorem{remarks}[subsection]{Remarks}

\numberwithin{equation}{subsection}

% Macros
%
\def\lim{\mathop{\mathrm{lim}}\nolimits}
\def\colim{\mathop{\mathrm{colim}}\nolimits}
\def\Spec{\mathop{\mathrm{Spec}}}
\def\Hom{\mathop{\mathrm{Hom}}\nolimits}
\def\Ext{\mathop{\mathrm{Ext}}\nolimits}
\def\SheafHom{\mathop{\mathcal{H}\!\mathit{om}}\nolimits}
\def\SheafExt{\mathop{\mathcal{E}\!\mathit{xt}}\nolimits}
\def\Sch{\mathit{Sch}}
\def\Mor{\mathop{\mathrm{Mor}}\nolimits}
\def\Ob{\mathop{\mathrm{Ob}}\nolimits}
\def\Sh{\mathop{\mathit{Sh}}\nolimits}
\def\NL{\mathop{N\!L}\nolimits}
\def\CH{\mathop{\mathrm{CH}}\nolimits}
\def\proetale{{pro\text{-}\acute{e}tale}}
\def\etale{{\acute{e}tale}}
\def\QCoh{\mathit{QCoh}}
\def\Ker{\mathop{\mathrm{Ker}}}
\def\Im{\mathop{\mathrm{Im}}}
\def\Coker{\mathop{\mathrm{Coker}}}
\def\Coim{\mathop{\mathrm{Coim}}}

% Boxtimes
%
\DeclareMathSymbol{\boxtimes}{\mathbin}{AMSa}{"02}

%
% Macros for moduli stacks/spaces
%
\def\QCohstack{\mathcal{QC}\!\mathit{oh}}
\def\Cohstack{\mathcal{C}\!\mathit{oh}}
\def\Spacesstack{\mathcal{S}\!\mathit{paces}}
\def\Quotfunctor{\mathrm{Quot}}
\def\Hilbfunctor{\mathrm{Hilb}}
\def\Curvesstack{\mathcal{C}\!\mathit{urves}}
\def\Polarizedstack{\mathcal{P}\!\mathit{olarized}}
\def\Complexesstack{\mathcal{C}\!\mathit{omplexes}}
% \Pic is the operator that assigns to X its picard group, usage \Pic(X)
% \Picardstack_{X/B} denotes the Picard stack of X over B
% \Picardfunctor_{X/B} denotes the Picard functor of X over B
\def\Pic{\mathop{\mathrm{Pic}}\nolimits}
\def\Picardstack{\mathcal{P}\!\mathit{ic}}
\def\Picardfunctor{\mathrm{Pic}}
\def\Deformationcategory{\mathcal{D}\!\mathit{ef}}
\usepackage{cleveref}
\crefname{originalreference}{original reference}{original references}
\Crefname{originalreference}{Original reference}{Original references}
\makeatletter
\newlabel{situation-setup}{{02QL}{source}{}{sourcepointer-1}{}}
\newlabel{situation-setup@cref}{{[originalreference][0][]02QL}{[1][0][]source}{}{}{}}
\makeatother
\crefname{originalreference}{original reference}{original references}
\Crefname{originalreference}{Original reference}{Original references}
\makeatletter
\newlabel{situation-setup-base-change}{{0FVG}{source}{}{sourcepointer-2}{}}
\newlabel{situation-setup-base-change@cref}{{[originalreference][0][]0FVG}{[1][0][]source}{}{}{}}
\makeatother
\crefname{originallemma}{original lemma}{original lemmas}
\Crefname{originallemma}{Original lemma}{Original lemmas}
\makeatletter
\newlabel{lemma-pullback-base-change}{{0FVJ}{source}{}{sourcepointer-3}{}}
\newlabel{lemma-pullback-base-change@cref}{{[originallemma][0][]0FVJ}{[1][0][]source}{}{}{}}
\makeatother
\crefname{originallemma}{original lemma}{original lemmas}
\Crefname{originallemma}{Original lemma}{Original lemmas}
\makeatletter
\newlabel{morphisms-lemma-base-change-flat}{{01U9}{source}{}{sourcepointer-4}{}}
\newlabel{morphisms-lemma-base-change-flat@cref}{{[originallemma][0][]01U9}{[1][0][]source}{}{}{}}
\makeatother
\crefname{originallemma}{original lemma}{original lemmas}
\Crefname{originallemma}{Original lemma}{Original lemmas}
\makeatletter
\newlabel{lemma-pullback-coherent-base-change}{{0FVI}{source}{}{sourcepointer-5}{}}
\newlabel{lemma-pullback-coherent-base-change@cref}{{[originallemma][0][]0FVI}{[1][0][]source}{}{}{}}
\makeatother
\title{Composition of two flat base changes is a flat base change and the pullbacks compose}
\author{The Stacks Project Authors}
\date{stacks-project:chow.tex:local-snapshot-20261012; GNU FDL 1.2 or later}
\begin{document}
\maketitle
\noindent\textit{Editorial note.} All delivered mathematical text is the author's own, copied from the pinned Stacks Project chapter file chow.tex; no statement or proof step was rewritten, reordered or removed.\par\medskip
\noindent\textit{Editorial note.} Preamble provenance: the pinned chapter file opens with an input of the shared preamble and carries no preamble text of its own, so the project's single Stacks-wide preamble staged at sources/stacks-chow/preamble.tex is excerpted verbatim as two spans, lines 14-18 (Xy-pic setup) and lines 154-217 (the native theorem declarations and mathematical macros). Nothing was added, removed or redefined; the wrapper class and the amsmath/amssymb/amsthm/graphicx base packages are supplied by the assembler.\par\medskip
\noindent\textit{Editorial note.} The counted statement and proof contain no cross reference, no citation and no included file, so the delivered body is self contained.\par\medskip
\noindent\textit{Editorial note.} Canonical location of the counted proposition: Stacks tag 0FVP (sources/stacks-chow/pinned-original-source.tex line 15438).\par\medskip
\noindent\textit{Allowed standard prerequisites (editorial).}\begin{itemize}
\item The base change situation and its shift constant. Same chapter, Situation labelled situation-setup-base-change, invoked in the statement. Fixes the hypotheses and the meaning of the shifts c and c'.
\item Lifting of generalizations along flat morphisms. Morphisms, Lemma labelled morphisms-lemma-base-change-flat, cited at proof lines 15470-15472. Makes the image of a generic point of a fibre component again a generic point of a fibre component.
\item Base change of cycles and its composition. Same chapter, Lemma labelled lemma-pullback-coherent-base-change, cited at proof lines 15478-15480. Gives (g')\textasciicircum{}*[Z']\_\{k+c\} = [Z'']\_\{k+c+c'\} and hence the composition on Chow groups.
\item The pullback map on Chow groups. Same chapter, Lemma labelled lemma-pullback-base-change, cited in the statement. Identifies the Chow level maps whose composition is asserted.
\end{itemize}\medskip
\ifdefined\hypertarget\hypertarget{sourcepointer-1}{}\fi
\noindent\textit{Source reference.} Original reference 02QL: sources/stacks-chow/pinned-original-source.tex, source lines 1372--1372. Same-chapter pointer into the author environment situation at source line 1372; invoked at proof line None of the counted proof. It is retained only as a navigation pointer; every load-bearing step is inside the delivered span.\par\medskip
\ifdefined\hypertarget\hypertarget{sourcepointer-2}{}\fi
\noindent\textit{Source reference.} Original reference 0FVG: sources/stacks-chow/pinned-original-source.tex, source lines 15123--15123. Same-chapter pointer into the author environment situation at source line 15123; invoked at proof line None of the counted proof. It is retained only as a navigation pointer; every load-bearing step is inside the delivered span.\par\medskip
\ifdefined\hypertarget\hypertarget{sourcepointer-3}{}\fi
\noindent\textit{Source reference.} Original lemma 0FVJ: sources/stacks-chow/pinned-original-source.tex, source lines 15255--15255. Same-chapter pointer into the author environment lemma at source line 15255; invoked at proof line None of the counted proof. It is retained only as a navigation pointer; every load-bearing step is inside the delivered span.\par\medskip
\ifdefined\hypertarget\hypertarget{sourcepointer-4}{}\fi
\noindent\textit{Source reference.} Original lemma 01U9: sources/stacks-chow/tags.txt, source lines 2402--2402. Cross-chapter standard prerequisite located in the project tag list staged in the local snapshot; invoked at proof line 15467 of the counted proof. It is retained only as a navigation pointer; every load-bearing step is inside the delivered span.\par\medskip
\ifdefined\hypertarget\hypertarget{sourcepointer-5}{}\fi
\noindent\textit{Source reference.} Original lemma 0FVI: sources/stacks-chow/pinned-original-source.tex, source lines 15193--15193. Same-chapter pointer into the author environment lemma at source line 15193; invoked at proof line 15478 of the counted proof. It is retained only as a navigation pointer; every load-bearing step is inside the delivered span.\par\medskip
\begin{lemma}
\label{lemma-compose-base-change}
Let $(S, \delta)$, $(S', \delta')$, $(S'', \delta'')$ be as in
Situation \ref{situation-setup}. Let $g : S' \to S$ and $g' : S'' \to S'$
be flat morphisms of schemes and let $c, c' \in \mathbf{Z}$
be integers such that $S, \delta, S', \delta', g, c$ and
$S', \delta', S'', g', c'$ are as in
Situation \ref{situation-setup-base-change}.
Let $X \to S$ be locally of finite type and denote $X' \to S'$
and $X'' \to S''$ the base changes by $S' \to S$ and $S'' \to S$.
Then
\begin{enumerate}
\item $S, \delta, S'', \delta'', g \circ g', c + c'$ is as in
Situation \ref{situation-setup-base-change},
\item the maps $g^* : Z_k(X) \to Z_{k + c}(X')$ and
$(g')^* : Z_{k + c}(X') \to Z_{k + c + c'}(X'')$ of
compose to give the map $(g \circ g')^* : Z_k(X) \to Z_{k + c + c'}(X'')$, and
\item the maps $g^* : \CH_k(X) \to \CH_{k + c}(X')$ and
$(g')^* : \CH_{k + c}(X') \to \CH_{k + c + c'}(X'')$ of
Lemma \ref{lemma-pullback-base-change}
compose to give the map $(g \circ g')^* : \CH_k(X) \to \CH_{k + c + c'}(X'')$
of Lemma \ref{lemma-pullback-base-change}.
\end{enumerate}
\end{lemma}

\begin{proof}
Let $s \in S$ and let $s'' \in S''$ be a generic point of an irreducible
component of $(g \circ g')^{-1}(\{s\})$. Set $s' = g'(s'')$.
Clearly, $s''$ is a generic point of an irreducible component of
$(g')^{-1}(\{s'\})$. Moreover, since $g'$ is flat and hence generalizations
lift along $g'$ (Morphisms, Lemma \ref{morphisms-lemma-base-change-flat})
we see that also $s'$ is a generic point of an irreducible component
of $g^{-1}(\{s\})$. Thus by assumption $\delta'(s') = \delta(s) + c$
and $\delta''(s'') = \delta'(s') + c'$. We conclude
$\delta''(s'') = \delta(s) + c + c'$ and the first part of the
statement is true.

\medskip\noindent
For the second part, let $Z \subset X$ be an integral closed subscheme
of $\delta$-dimension $k$. Denote $Z' \subset X'$ and $Z'' \subset X''$
the base changes. By definition we have $g^*[Z] = [Z']_{k + c}$.
By Lemma \ref{lemma-pullback-coherent-base-change} we have
$(g')^*[Z']_{k + c} = [Z'']_{k + c + c'}$. This proves the final statement.
\end{proof}

\end{document}
